\section{Implementation}
\label{sec:impl}

An implementation keeps the guarantees of our methods if it realizes the loads, stores, CASes and fences of the model of Section~\ref{sec:bg-model} as the model describes them: each atomic, in the order that the model keeps, and in a bounded number of steps of its thread. This section describes how our implementations, written in C++ and built with GCC for x86-64 (Section~\ref{sec:exp-setup}), do so; Table~\ref{tab:params} lists their parameters.

\begin{table}
    \caption{The parameters of our implementations, with the sections that describe them.}
    \label{tab:params}
    \small
    \begin{tabular}{l r r}
        \toprule
        Parameter & Value & Section \\
        \midrule
        $\alpha$ and $\beta$ of the direction rule & 15, 18 & \ref{sec:bg-bfs} \\
        Vertices per bottom-up block, $W$ & 1,024 & \ref{sec:do-bu} \\
        Neighbors per chunk, $g$ & 256 & \ref{sec:impl-chunks} \\
        Longest list a helper takes in order & 3 & \ref{sec:td-contention} \\
        Positions of a bucket's first array, $c_0$ & 32 & \ref{sec:td-buckets} \\
        Bytes per arena chunk, at least & 4~MiB & \ref{sec:impl} \\
        Frontier codes & 1 to 255 & \ref{sec:impl-codes} \\
        Random generator of a thread & \texttt{mt19937} & \ref{sec:impl} \\
        \bottomrule
    \end{tabular}
\end{table}

Helpers read buckets, progress fields and flags while their owners write them, and a thread that has fallen behind writes distances, codes and stamps that others are reading; with ordinary accesses, these would be data races, whose behavior C++ leaves undefined, and which let a compiler split an access, or load a location again instead of keeping a value that it loaded~\cite{boehm_cpp, boehm_benign}. Every access to a location that another thread may access at the same time, with one of the two a write, is therefore atomic: a \texttt{std::atomic} field (Table~\ref{tab:orders}, last row), or an ordinary object accessed with \texttt{std::atomic\_ref}, so that the arrays, and the code that sets them up and resets them while no other thread accesses them, stay ordinary. Each such object is a byte, a 32-bit or 64-bit integer, or a pointer, at the natural alignment that \texttt{atomic\_ref} requires for it on x86-64, and lives at least until every call of the search has returned; the compiler neither splits these accesses nor repeats or invents them. Table~\ref{tab:orders} gives the order of each kind of shared access: a \emph{release} store, which no earlier access of its thread may follow, publishes work or data, and an \emph{acquire} load that reads its value orders after it whatever the storing thread did before the store, ordinary accesses included; a \emph{relaxed} access orders nothing else.

\begin{table}
    \caption{The order of each kind of shared access in our implementations.}
    \label{tab:orders}
    \footnotesize
    \setlength{\tabcolsep}{3pt}
    \begin{tabular}{@{}>{\raggedright\arraybackslash}p{0.18\columnwidth} >{\raggedright\arraybackslash}p{0.57\columnwidth} >{\raggedright\arraybackslash}p{0.19\columnwidth}@{}}
        \toprule
        Role & Accesses & Order \\
        \midrule
        Publish work or data & a flag, the certificate of a chunk, a stamp or a progress field, after the work that it covers (Section~\ref{sec:td-order}); the code and then the distance of a vertex, after the vertex is appended (Section~\ref{sec:impl-codes}); in \textsc{Append}, the new size, after the vertex, and the address of a new array, after the copy (Algorithm~\ref{alg:bucket}); in the variant, the tag of a bucket, after the progress of its thread is reset and the bucket emptied (Algorithm~\ref{alg:mem}) & release store \\
        \addlinespace
        Accept a publication & the loads of these flags, certificates, stamps and progress fields, and of distances, before the work that they may let a thread skip; of the size and the array of another thread's bucket, before its vertices; in the variant, of the tag of a bucket, before the size or the progress that the tag guards (Algorithms~\ref{alg:mem} and~\ref{alg:mem-expand}) & acquire load \\
        \addlinespace
        Ordered already, or advisory & the accesses of a thread to its own bucket and fields; loads that an earlier acquire load or CAS of the thread already orders, such as those of the vertices of a bucket after its size, and the frontier tests of a bottom-up level; stores that a later release store publishes, such as the pass's codes; advisory loads, such as those of the sizes and edge counts that choose a direction & relaxed \\
        \addlinespace
        Publish a plan, node, record or discovery & the fields $\mathit{planner}$ and $\mathit{next}$, and the slots of the variant, which are \texttt{std::atomic} objects, and, with a CAS only, the distance in \textsc{Discover} (Algorithm~\ref{alg:wf-do}) & sequentially consistent CAS and load \\
        \bottomrule
    \end{tabular}
\end{table}

\paragraph{Buckets and their arrays.} An owner whose bucket is full copies its vertices into a larger array with ordinary stores before the release store of its address (Algorithm~\ref{alg:bucket}), so the copy happens before every read of the new array by a thread that has loaded the address. The copy reads the old array, and the size and address of the bucket's array, with relaxed loads, since helpers may be reading them at the same time, and C++ asks that an object referenced by an \texttt{atomic\_ref} be accessed only through one while that reference exists. The other ordinary accesses are to memory that only one thread uses during a search, such as the capacity of a bucket; to memory that no thread writes during a search, such as the graph; and to memory that a thread writes before publishing it with a CAS, which others read only after reading what the CAS wrote: a node before it is linked (Lemma~\ref{lem:list}), a proposal before its claim (Lemma~\ref{lem:plans}), and, in the variant, a record before it is installed (Lemma~\ref{lem:mem-slots}).

\paragraph{Arenas.} Each thread has two arenas, one for the arrays of its buckets and one for the nodes that it adds, or, in the variant, the records that it proposes. An arena hands out blocks from chunks of at least 4~MiB, which it keeps from one search to the next, and asks the system allocator for a new chunk only when no later chunk has room for the block. Keeping the arrays apart means that their memory only ever holds integers: allocating a chunk, with \texttt{operator new} or \texttt{calloc}, creates them implicitly~\cite{cppstd}, while memory that had held a node would need new integers created explicitly before an array could use it. In the variant, a thread that reads a reused bucket can read positions that no thread has written in the search (Section~\ref{sec:mem-reuse}), and C++ forbids reading the indeterminate value of memory never written, so its arenas take zeroed memory from \texttt{calloc}: every position then holds 0 or a vertex written earlier, which \textsc{Valid} discards unless it is of the reader's level, and the zeroing, like the copy, happens before every read of an array published in that memory. The system allocator may take locks; the method assumes only that allocation takes a bounded number of steps (Section~\ref{sec:bg-model}), which an implementation can meet by reserving, before any call begins, the memory that a search can allocate, bounded since a thread appends each vertex at most once (Lemma~\ref{lem:once}) and a search has at most $n+1$ levels (Section~\ref{sec:mem-memory}). Our arenas grow on demand instead, mostly in the first search of a thread count, and are reused after it.

\paragraph{Instructions and steps.} On x86-64, GCC compiles every relaxed, release or acquire load or store to a single \texttt{mov} or \texttt{movzx} instruction, which TSO keeps in program order with the other accesses of its thread, except a store and a later load, an order that neither release nor acquire asks for; these orders thus constrain only the compiler, and every thread makes its loads and stores in the order that the algorithms require. Every CAS is a sequentially consistent \texttt{compare\_exchange\_strong}, which GCC compiles to a \texttt{lock cmpxchg}, and the fence of Algorithm~\ref{alg:wf-td} a sequentially consistent \texttt{atomic\_thread\_fence}, a \texttt{lock or} of zero into the top of the stack; a locked instruction takes effect only once the earlier stores of its thread are visible~\cite{intelmanual}, and GCC moves no memory access across either, so both behave as Section~\ref{sec:bg-model} assumes. Every atomic type that we use, of one, four or eight bytes, is always lock-free on x86-64 (\texttt{is\_always\_lock\_free}), and each of our atomic operations is a single instruction, so each takes a bounded number of steps, as the proof of wait-freedom counts it. On ARM and POWER, GCC emits instructions of their own for release and acquire, but C++ promises less with these orders than TSO, since a release store need not become visible to all other threads at the same time, so our proofs do not cover such processors as they stand.

\paragraph{Random orders.} Every random order that our implementations draw, of the entries of a node, the positions of a bucket, the neighbors of a vertex, the blocks of a thread and the chunks of a vertex (Section~\ref{sec:impl-chunks}), is an \emph{affine permutation} of the $k$ items it orders, which visits item $(a + i b) \bmod k$ in step $i$, for a random $a$ and a random stride $b$ coprime to $k$; it takes two random numbers and no memory, while a uniformly random order of $k$ items takes $k$ random numbers and memory for all of them. An order of $k \leq 1$ items is the identity, and a thread skips an empty range of positions without drawing one. A thread draws its random numbers from an \texttt{mt19937} of its own, which refills its state once every 624 draws, and brings each into its range with a remainder rather than by rejection sampling, at the cost of a slight bias, so that a draw takes a bounded number of steps; the search for a stride coprime to the number of items takes at most $2\sqrt{k}$ tries (Lemma~\ref{lem:affine}). Threads repeat nearly as little work with affine orders as with uniform ones (Section~\ref{sec:exp-work}), although Theorem~\ref{thm:records} assumes uniform ones. Over two sockets, a thread visits the entries of threads on its own NUMA node first, in one random order, and then those on the other node, in another.

\paragraph{The reset and the result.} Our implementation resets the arrays of a search in parallel, and no thread starts to search before every share is reset, since a thread still resetting could overwrite a distance that another has already set. This is the one step of a search at which our threads wait for each other; it lies outside the method, which starts from arrays that hold their initial values (Section~\ref{sec:bg-method}), and Section~\ref{sec:exp-speed} measures it apart from the search. Calls that are still running when the first returns may write the distances again, with the values already there, so a consumer that reads them before every call has returned must read them with atomic loads, or first copy them with atomic loads into memory of its own.

\subsection{Frontier Codes}
\label{sec:impl-codes}

The test of a bottom-up level, $\mathit{dist}[u] = d$ (Section~\ref{sec:do-bu}), is the most frequent memory access of the bottom-up step, and $\mathit{dist}$ has four bytes per vertex. Our implementations test a denser array: one byte per vertex, $\mathit{code}[v]$, cleared with the other arrays before every search, and a \emph{code} $c_d$ for each level $d$ that a bottom-up level tests, so that the test becomes $\mathit{code}[u] = c_d$. Codes are handed out by the plans. Each plan carries $c_d$, or 0 if level $d$ tests distances, and $\ell_d$, the largest code handed out in the search so far; the thread that proposes the plan computes both from the plan of the level before, and the plan before level 0 has both 0. A top-down level gets $c_d = 0$ and keeps $\ell_d = \ell_{d-1}$. A bottom-up level takes $k = c_{d-1} + 1$ if level $d-1$ was bottom-up with $c_{d-1} > 0$, and $k = \ell_{d-1} + 1$ otherwise. Codes fit in a byte: if $k \leq 255$, then $c_d = k$ and $\ell_d = \min(k + 1, 255)$, and otherwise $c_d = 0$ and $\ell_d = \ell_{d-1}$. The discoveries of a bottom-up level $d$ write $r_d$ after appending a vertex and before writing its distance, where $r_d = c_d + 1$ if $0 < c_d < 255$, the code that the plan reserved, and $r_d = 0$ otherwise. A level with $c_d = 0$ tests $\mathit{dist}[u] = d$; a test of code 0 would be wrong, since an undiscovered vertex has code 0. Once a bottom-up level gets the code 0, $\ell = 255$, and every later level tests distances too. We say that $c_d$, if positive, belongs to level $d$, and $r_d$, if positive, to level $d+1$, which tests it if it is bottom-up, since then $c_{d+1} = r_d$, and leaves it unused otherwise. Each positive code belongs to only one level: $\ell$ never decreases, the codes that the plan of level $d$ hands out, $c_d$ and $r_d$, are at most $\ell_d$, and a code $k = \ell_{d-1} + 1$ is larger than every code handed out before it. A bottom-up level with $c_d > 0$ that follows a top-down one, whose appends write no codes, starts with a pass that writes $c_d$ for the vertices in the buckets of $N_d$. The pass is shared as an expansion is: a thread writes the codes of its own bucket in order, publishing how far it has got in a field $\mathit{codesLast}$ of its entry, then writes the rest of every other bucket whose flag $\mathit{codesDone}$ is not set, and sets that flag once the bucket's codes are all written. In the variant, a helper writes codes only for vertices that pass \textsc{Valid}. A thread starts testing only after the pass. Top-down levels write no codes, since threads appending nearby vertices would contend for the cache lines of $\mathit{code}$.

With codes, a write can also be unsound by writing a wrong code, so we add a third safety property to those of Section~\ref{sec:pf-safety}:
\begin{itemize}
    \item[(S3)] every positive value written to $\mathit{code}[v]$ belongs to level $\delta(v)$.
\end{itemize}

\begin{lemma}[Frontier codes]
\label{lem:codes}
Let level $d$ be bottom-up, with $c_d > 0$. (a) If every write so far is sound and a thread reads $\mathit{code}[u] = c_d$, then $\delta(u) = d$. (b) If every write so far is sound, and (R1) and (R2) hold for $N_{d-1}$ and $N_d$ from the times their plans are claimed, then when a thread in level $d$ tests a vertex $u \in L_d$, it finds $\mathit{code}[u] = c_d$.
\end{lemma}
\begin{proof}
(a) The array is cleared before the search, so the thread read a value that a write of the search stored. By (S3), that value belongs to level $\delta(u)$, and $c_d$ belongs to level $d$ and to no other, so $\delta(u) = d$.
(b) Every write to $\mathit{code}[u]$ writes $c_d$. By (S3), a positive value written there belongs to level $d$, so it is $c_d$ or $r_{d-1}$, and $r_{d-1}$, if positive, is $c_{d-1} + 1 = c_d$. A write of 0 would be a discovery of a bottom-up level $d-1$ with $c_{d-1} = 0$ or $c_{d-1} = 255$, after which $c_d = 0$. If level $d-1$ was bottom-up, $c_{d-1} > 0$, since $c_d > 0$; every vertex of $L_d$ had been discovered in level $d-1$ by the time the first thread finished $N_{d-1}$ (Lemma~\ref{lem:bu-complete}), which is before $\tau_d$, and each discovery wrote $r_{d-1} = c_d$ before the distance. If level $d-1$ was top-down, the thread made the pass before testing, and the pass gave every vertex of $L_d$ its code: $\mathit{codesLast}$ and $\mathit{codesDone}$ are written after the codes they certify, as $\mathit{last}$ and $\mathit{done}$ are after the expansions they certify (Lemma~\ref{lem:certs}(b)), and every vertex of $L_d$ is in a reference prefix by (R2), so the argument of Lemma~\ref{lem:td-complete} applies with codes in place of expansions. Under TSO, every code is stored before the distance or the field that certifies it, and the argument of Section~\ref{sec:pf-tso} applies.
\end{proof}

The proof of Lemma~\ref{lem:ready} carries over to the method with codes, with (S3) among the properties of a sound write. Let $w$ be the first write that is not sound. In the induction on $d$, Lemma~\ref{lem:bu-complete} holds for a level that tests codes, since the test $\mathit{code}[u] = c_d$ holds for every vertex $u$ of $L_d$ by part~(b) of the lemma, in place of (R1); (R1) and (R2) thus hold for every level whose plan is claimed before $w$, from then up to $w$. Then $w$ cannot be a discovery in a bottom-up level $e$ that tests codes, nor its append: the neighbor $u$ that the thread found with $\mathit{code}[u] = c_e$ has $\delta(u) = e$ by part~(a), so the proof of Lemma~\ref{lem:discover} applies. Nor can $w$ be the code $r_e$ that such a discovery writes, since the discovery is sound, so $\delta(v) = e + 1$, and $r_e$ belongs to level $e + 1$. Nor can $w$ be a code that the pass of a level $e$ writes, since $w$ writes $c_e$, which belongs to level $e$, for a vertex of a bucket of $N_e$, which is in $L_e$ by (S2), or, in the variant, passes \textsc{Valid} and is in $L_e$ by Lemma~\ref{lem:mem-reads}. All other writes are those of Section~\ref{sec:proof}. Every write is therefore sound, and Theorems~\ref{thm:correct} and~\ref{thm:waitfree} hold for the method with codes; the pass adds a constant number of steps per position of a bucket to the count of Theorem~\ref{thm:waitfree}.

\subsection{Chunk Certificates}
\label{sec:impl-chunks}

In the method, a thread that enters the expansion of a vertex examines all of its neighbors (Section~\ref{sec:td-contention}), so a schedule can make every thread examine again every neighbor of a vertex of high degree (Proposition~\ref{prop:granularity}). Our implementations therefore certify the neighbors of a vertex of more than $g$ neighbors in \emph{chunks} of $g$ consecutive neighbors, the last of which may be shorter, as part~(b) of the proposition describes; we use $g = 256$. Chunk $c$ of a vertex $u$ has a \emph{certificate} $\mathit{cert}[u][c]$, which a thread sets to $\nu$, the number of the current search, once it has examined every neighbor in the chunk, and other threads skip a chunk whose certificate holds $\nu$. A thread that expands $u$ goes through its chunks, scanning each one that is not certified and certifying it afterwards, and sets $\mathit{expanded}[u]$ once it has been through every chunk. A thread that stops in the middle of a chunk thus leaves uncertified only the at most $g$ neighbors of that chunk that it has examined; the chunks that it has not started stay whole for the others, which do not scan again the chunks certified before it. An owner takes the $K$ chunks of a vertex in order from chunk $\lfloor tK/T \rfloor$, so that owners that hold copies of the same vertex start apart, and a helper takes them in a random order of its own, and the neighbors of each chunk in another. A vertex of at most $g$ neighbors is expanded as in Algorithm~\ref{alg:wf-td}, since $\mathit{expanded}[u]$ certifies the same work as a single chunk would.

The certificates are one array of 32-bit words, with the certificates of a vertex $u$ at $\lfloor \mathit{offsets}[u] / g \rfloor + u$ onward, $n + m/g + 1$ words in all, about as much memory as the distances take; a vertex of at most $g$ neighbors leaves its word unused. Offsets into the adjacency array, like edge counts, are 64-bit integers, since that of Twitter has about 2.9 billion entries; vertices are 32-bit integers. The words of two vertices never overlap: with $o = \mathit{offsets}[u]$ and $q = \deg(u)$, the words of $u + 1$ start $\lfloor (o + q)/g \rfloor - \lfloor o/g \rfloor + 1 \geq \lfloor q/g \rfloor + 1 \geq \lceil q/g \rceil$ words after those of $u$, and $u$ has $\lceil q/g \rceil$ chunks. A search writes its number into them, and the numbers grow from search to search, so a certificate left by an earlier search holds a smaller number, and the array is cleared only before the first search of each thread count, whose numbers start again at 1. The numbers are positive 32-bit integers, since our implementation uses negative ones to stop its threads, so this holds for $2^{31}-1$ searches; a program that runs more would have to clear the array again before the numbers wrap around, once every call of the last search has returned.

\begin{lemma}[Chunk certificates]
\label{lem:chunks}
When a thread certifies a chunk of a vertex $u$ in the search, every neighbor $v$ of $u$ in the chunk with $\delta(v) = \delta(u) + 1$ has been appended to a bucket of $N_{\delta(v)}$ and has $\mathit{dist}[v] = \delta(v)$, and with chunks, Lemma~\ref{lem:certs}(a) still holds.
\end{lemma}
\begin{proof}
A thread sets $\mathit{cert}[u][c]$ only after examining every neighbor of $u$ in chunk $c$, and for each such neighbor $v$, as in the proof of Lemma~\ref{lem:certs}(a), it either appended $v$ before writing its distance, or found the distance set by a thread that had appended $v$ first. Certificates written in earlier searches hold smaller numbers than $\nu$, so a certificate equal to $\nu$ was written in this search, by a thread that had scanned the chunk first. A thread sets $\mathit{expanded}[u]$ only after going through every chunk of $u$, each of which it scanned or found certified with $\nu$, so by then every neighbor $v$ of $u$ with $\delta(v) = \delta(u) + 1$ has been appended and has its distance, which is Lemma~\ref{lem:certs}(a).
\end{proof}

With chunks, a call of \textsc{Expand} expands $u$ completely if, for every chunk of $u$, it examines every neighbor in the chunk or reads its certificate equal to $\nu$, and then sets $\mathit{expanded}[u]$. The proofs of Section~\ref{sec:proof} use complete expansions only through Lemma~\ref{lem:certs}(a) and the fact that a complete expansion sets $\mathit{expanded}[u]$, so by the lemma, Theorems~\ref{thm:correct} and~\ref{thm:waitfree} hold for the implementation with chunks. A call reads each certificate at most once, which adds at most $n + m/g$ steps to the bound of Theorem~\ref{thm:waitfree}, and a top-down level $d$ takes $O(T(Q_d + \mathrm{vol}(L_d) + C_d) + T^2)$ steps in Theorem~\ref{thm:work-td}, where $C_d \leq |L_d| + \mathrm{vol}(L_d)/g$ counts the chunks of vertices of $L_d$.

The argument of Theorem~\ref{thm:records} applies to chunks taken in uniformly random orders as well, if there are at most $n$ threads: the threads that expand a vertex scan $O((k + K)\log(n+m))$ of its $K$ chunks, where $k$ is the number of times threads enter its expansion, instead of up to $kK$. Combined with the theorem, this bounds the neighbors that a level examines.

\begin{corollary}[Neighbors examined]
\label{cor:neighbors}
Suppose that $n \geq 2$ and $T \leq n$; that the assumption of Theorem~\ref{thm:records} on memory holds for the certificates of chunks as well as for the flags $\mathit{expanded}$, so that memory is sequentially consistent or a thread executes a fence after every certificate; that helpers take positions as the theorem assumes; and that before the search, each helper draws, for each vertex, an order of the vertex's chunks uniformly at random, independently of all other draws, and takes the chunks of every vertex that it expands in that order. Let $\lambda = \log(n + m + 2)$. Then for every constant $c$, with probability at least $1 - n^{-c}$ and for every schedule, the helpers of a top-down level $d$ examine $O\bigl((\mathrm{vol}(L_d) + g Q_d + g T^2)\lambda^2\bigr)$ neighbors; if each owner also draws its orders of chunks in this way, then so do all threads of the level together.
\end{corollary}
\begin{proof}
Let $k_u$ be the number of times threads enter the expansion of $u \in L_d$, counting only helpers for the first claim. By Theorem~\ref{thm:records}, and since an owner expands each of its positions at most once, $\sum_{u} k_u = O((T^2 + Q_d)\log n)$. The proof of Theorem~\ref{thm:records} applies to the $K_u = \lceil \deg(u)/g \rceil$ chunks of $u$ and the expansions of $u$, with Lemma~\ref{lem:records} for certificates, which needs the assumption on memory, and with the union bound taken also over which threads enter them. There are at most $T^{k_u} \leq n^{k_u}$ such choices, since $T \leq n$, and at most $K_u! \leq m^{K_u}$ orders in which the chunks are certified, since $K_u \leq \deg(u) \leq m$; the degree has no bound in terms of $n$ alone, since our graphs keep repeated edges (Section~\ref{sec:exp-setup}). So the union bound adds $O((k_u + K_u)\lambda)$ to the exponent, and with $z = a(k_u + K_u)\lambda$ for $a$ large enough, it bounds the chunks of $u$ that these expansions scan by $O((k_u + K_u)\lambda)$, whatever order the other threads take them in. A chunk holds at most $g$ neighbors, so these expansions examine $O(g\lambda \sum_u (k_u + K_u)) = O\bigl(g\lambda\,((T^2 + Q_d)\log n + |L_d| + \mathrm{vol}(L_d)/g)\bigr)$ neighbors, which is the bound, since $|L_d| \leq Q_d$ and $\log n \leq \lambda$.
\end{proof}

Our owners take the chunks of a vertex in order, from a start that depends on the thread, which keeps their scans sequential. For a hypothetical implementation whose helpers draw independent uniform orders, and whose marks are ordered as Theorem~\ref{thm:records} requires, the corollary would still bound the neighbors that helpers examine, but not those that owners examine, since a vertex can be in the buckets of several owners, and with fixed orders the argument does not apply to them. Our helpers draw affine orders, and our marks are stores without a fence, so the corollary does not cover our implementation, whose repeated work Section~\ref{sec:exp-work} measures instead, and Section~\ref{sec:exp-hubs} what chunks save.
